\section{Associativity by counting filtrations}\label{hall:chapter}

An extension records an attachment. We can also ask how many attachments
of one type occur inside a fixed middle object. This question leads to a
new multiplication. Its associativity will follow from the correspondence
between submodules of a quotient and submodules containing its kernel.
The quotient calculation that began the exposition will thus return as
the mechanism behind an algebra law.

\subsection{A setting where the counts are finite}

Fix a finite field $\kk$ with $q$ elements. We work with finite-dimensional
$\kk[t]$-modules on which $t$ acts nilpotently. Equivalently, our objects
are pairs $(V,N)$ with $N^s=0$ for some $s$, and our maps are linear maps
commuting with the operators. A submodule is exactly an $N$-invariant
subspace. Throughout this chapter an isomorphism must respect $N$.

Submodules and quotients stay in this class because the same power of
$N$ is zero on them. Extensions stay in the class too. Indeed, in an
exact sequence
\[
0\longrightarrow B\longrightarrow E\longrightarrow A\longrightarrow0,
\]
if $t^aA=0$ and $t^bB=0$, then $t^aE\subseteq B$ and hence
$t^{a+b}E=0$. Also $\dim E=\dim A+\dim B$, by the proved dimension
formula for an exact sequence. These are closure checks, not extra
assumptions about the category.

For isomorphism classes $A,B,E$, define the nonnegative integer
\begin{equation}\label{hall:count}
g_{A,B}^{E}
=\#\{U\subseteq E:\ U\text{ is a submodule},\ U\cong B,\ E/U\cong A\}.
\end{equation}
The order matters: $B$ is the lower layer and $A$ is the quotient above it.
This counts actual submodules of $E$. It does not count chosen
isomorphisms $B\to U$, chosen quotient identifications, or marked
extension classes. An isomorphism $E\to E'$ gives a bijection on these
sets by sending $U$ to its image, so the count depends only on the three
isomorphism classes.

The set in \eqref{hall:count} is finite because $E$ has finitely many
elements. For fixed $A,B$, only finitely many classes $E$ can occur.
They have the fixed dimension $d=\dim A+\dim B$. Choosing a basis
identifies such an operator with a $d$ by $d$ matrix over $\kk$, and
there are $q^{d^2}$ matrices altogether. Each isomorphism class has a
representative among them. This argument establishes finiteness without
requiring a classification of nilpotent operators.

Let $\mathcal H$ be the rational vector space with one basis symbol
$u_M$ for each isomorphism class in our collection. Explicitly, its
vectors are finite rational linear combinations of these symbols.
Define a bilinear multiplication by
\begin{equation}\label{hall:product}
u_Au_B=\sum_{[E]}g_{A,B}^{E}u_E.
\end{equation}
The preceding finiteness checks make this an element of $\mathcal H$.
It is the untwisted Hall multiplication for this collection of modules.
We now prove its algebra laws.

\subsection{One flag, two intermediate objects}

\begin{theorem}[Associativity from a flag bijection]\label{hall:associativity}
The multiplication \eqref{hall:product} is associative, with identity $u_0$.
\end{theorem}

\begin{proof}
Fix $A,B,C,E$. Count the set of flags
\begin{equation}\label{hall:flag}
0\subseteq U\subseteq V\subseteq E,
\qquad U\cong C,\quad V/U\cong B,\quad E/V\cong A.
\end{equation}
Inclusions may be equal when one of the specified factors is zero.

First fix $U$. If $E/U\cong D$, there are $g_{D,C}^{E}$ choices of $U$.
For each such $U$, the possible submodules $V/U\subseteq E/U$ with
the required factors are counted by $g_{A,B}^{D}$. The quotient
correspondence gives a bijection
\[
\{V:U\subseteq V\subseteq E\}
\longleftrightarrow
\{W:W\subseteq E/U\},
\quad V\longmapsto V/U,\quad W\longmapsto\pi^{-1}(W).
\]
It also gives $(E/U)/(V/U)\cong E/V$, by the map
$(e+U)+(V/U)\mapsto e+V$. Thus the number of flags is
\[
\sum_{[D]}g_{A,B}^{D}g_{D,C}^{E}.
\]
This is the coefficient of $u_E$ in $(u_Au_B)u_C$.

Alternatively fix $V$. If $V\cong D$, there are $g_{A,D}^{E}$ choices
of $V$, and then $g_{B,C}^{D}$ choices of $U\subseteq V$.
The same flags are therefore counted by
\[
\sum_{[D]}g_{B,C}^{D}g_{A,D}^{E},
\]
the coefficient of $u_E$ in $u_A(u_Bu_C)$. No choice of an isomorphism
to a representative $D$ affects either count: every such isomorphism
induces a bijection of the relevant submodule sets. Both counting
procedures identify their data with precisely the flags in
\eqref{hall:flag}. Equality of all coefficients proves associativity on
basis vectors, and bilinearity proves it on $\mathcal H$.

For the unit, the only submodule of $E$ isomorphic to zero is $0$,
and the only submodule with zero quotient is $E$. Consequently
$g_{A,0}^{E}=g_{0,A}^{E}$ is $1$ if $E\cong A$ and $0$ otherwise.
Hence $u_Au_0=u_0u_A=u_A$.
\end{proof}

The two readings of the flag use different intermediate objects:
\[
\begin{tikzcd}[column sep=large,row sep=large]
0\arrow[r]&U\arrow[r,hook]\arrow[d,equal]&V\arrow[r,two heads]\arrow[d,hook]
  &V/U\arrow[r]\arrow[d,hook]&0\\
0\arrow[r]&U\arrow[r,hook]&E\arrow[r,two heads]&E/U\arrow[r]&0.
\end{tikzcd}
\]
Both rows are exact. The right vertical map sends $v+U$ to $v+U$;
it is injective since its zero class means $v\in U$. The square on the
right commutes on each representative $v$. The quotient of that vertical
inclusion is $E/V$. Thus the intermediate object is either $V$, carrying
the lower two layers, or $E/U$, carrying the upper two layers.

At the beginning of the exposition, associativity of the group operation
was an axiom permitting us to compose translations. Here associativity
of a newly defined operation is a theorem. The theorem is proved by a
bijection: the two bracketings organize exactly the same data. This is
a structural explanation with an explicit inverse correspondence, not
an appeal to a resemblance between diagrams.

\subsection{The first product, computed completely}

Write $S=\kk[t]/(t)$ and $J_r=\kk[t]/(t^r)$. Thus $J_r$ is the single
nilpotent chain of length $r$ constructed in the preceding chapter.
There are two isomorphism classes of two-dimensional nilpotent modules:
$S\oplus S$ and $J_2$. One can see this directly. If $N=0$, the first
description holds. If $N\neq0$, a two-dimensional nilpotent operator has
a chain $Nv\neq0$, $N^2v=0$, and $(Nv,v)$ is a basis. For completeness,
one may obtain such a chain by taking the last two nonzero vectors in
any maximal chain; independence of a chain of length three would
contradict dimension two. The chain basis gives $J_2$.

For $S\oplus S$, every line is a submodule of type $S$ with quotient
$S$. There are $q+1$ lines: the $q$ lines spanned by $(1,a)$ for
$a\in\kk$, and the line spanned by $(0,1)$. In $J_2$, every invariant
line has zero restricted operator, since a nilpotent scalar is zero in
a field. It must therefore lie in $\Ker N$, which is itself one line.
There is exactly one choice. Hence
\begin{equation}\label{hall:square}
u_S^2=(q+1)u_{S\oplus S}+u_{J_2}.
\end{equation}

The nonsplit middle object appears once in this product, although the
marked extension calculation gave $q-1$ nonzero classes when the
endpoints are identified with $S$. The reason is the definition of the
count. Equation~\eqref{hall:square} counts submodules of a fixed middle
object. Marked extension equivalence keeps endpoint identifications;
our Hall basis keeps only the middle object's isomorphism class.
Neither convention can silently replace the other.

\subsection{Three layers and a second count}

The preceding reconstruction theorem gives exactly three nilpotent
types in dimension three:
\[
S^{\oplus3},\qquad J_2\oplus S,\qquad J_3.
\]
The coefficient of $u_E$ in $u_S^3$ counts full invariant flags
$0\subset U\subset V\subset E$, with $\dim U=1$ and $\dim V=2$.
Every factor is $S$, because a one-dimensional nilpotent operator is zero.

For $E=S^{\oplus3}$ all flags are invariant. There are
$(q^3-1)/(q-1)=q^2+q+1$ choices of $U$: the nonzero vectors partition
into sets of $q-1$ vectors spanning the same line. For each $U$, the
planes containing it correspond to the $q+1$ lines in $E/U$.
This gives $(q^2+q+1)(q+1)$ flags.

For $E=J_2\oplus S$, put $L=\im N$ and $K=\Ker N$. Then
$L\subset K$, $\dim L=1$, and $\dim K=2$. An invariant line $U$ lies
in $K$, giving $q+1$ choices. An invariant plane $V$ contains $L$:
the induced operator on the one-dimensional quotient $E/V$ is
nilpotent and hence zero, so $N(E)\subseteq V$. Conversely, every plane
containing $L$ is invariant. If $U=L$, there are $q+1$ possible planes
through it. For each of the other $q$ lines $U\subset K$, the only
plane containing both $U$ and $L$ is $K$. The count is $2q+1$.

For $J_3$, an invariant line must equal $\Ker N$, and an invariant
plane must equal $\im N$ by the same quotient argument and their
dimensions. These two subspaces form a flag, so the count is $1$.
We obtain
\begin{equation}\label{hall:cube}
u_S^3=(q+1)(q^2+q+1)u_{S^{\oplus3}}
       +(2q+1)u_{J_2\oplus S}+u_{J_3}.
\end{equation}
Associativity says that either order of the two successive Hall
multiplications gives this count. The flag description explains why
before any structure constants are expanded.

\subsection{What the multiplication remembers}

Every nonzero coefficient in $u_Au_B$ comes from a short exact sequence,
so it satisfies
\[
[E]=[A]+[B]\quad\text{in }K_0.
\]
This makes Hall multiplication graded by the additive invariant: the
degree of $u_M$ is $[M]$. In our nilpotent collection the only simple
module is $S$. Indeed, if a nonzero nilpotent module is simple, its
nonzero kernel of $t$ is the whole module; then every vector line is a
submodule, forcing dimension one. Writing $\mathcal N_q$ for this
nilpotent collection, we have $K_0(\mathcal N_q)\cong\mathbb Z$, and
this grading is ordinary dimension. Here $K_0(\mathcal N_q)$ is formed
using generators and exact-sequence relations inside $\mathcal N_q$,
as in Section~\ref{k:section}; the closure checks at the beginning of
this chapter make that construction applicable.

Within a single degree, the Hall algebra retains different basis vectors
for different isomorphism classes. The terms in \eqref{hall:square} all
have degree two, but the multiplication distinguishes the split object
from the chain and records their different numbers of subobjects.
This statement is a grading statement; it does not claim that $K_0$ is
a quotient algebra of $\mathcal H$.

We have reached three distinct forms of information. Composition
multiplicities retain what is additive across exact sequences.
Marked extensions retain a specified two-layer attachment. Hall
coefficients retain counts of subobjects and their quotients, and their
associativity packages the compatibility of longer flags. In the
nilpotent setting, the spaces of maps studied earlier recover the whole
isomorphism class. Each construction answers a precise question about
what the preceding description left unresolved.

Hall algebras can be developed much further, including settings where
geometric invariants replace finite counts. That further theory requires
additional constructions. The result established here is the finite
counting algebra and its associativity, proved directly from flags.
